In general, the cartesian square of~3.7 defines a morphism of \(X\)-functors:
\[
\xymatrix{
\mathrm{T}_{X/S}(\mathcal{M})\ar[rr]\ar[rd]
& & \mathrm{T}_{X'/S}(\mathcal{M})\times_{X'}X\ar[ld]\\
& X. &
}
\]

\begin{proposition}{3.7. bis}\label{II.3.7.bis}
Let \(f\colon X\to X'\) be an \(S\)-morphism; if \(X\) and \(X'\) satisfy
\textup{(E)} with respect to~\(S\), then
\[
\mathrm{T}_{X/S}(\mathcal{M})\xrightarrow{\mathrm{T}(f)}
\mathrm{T}_{X'/S}(\mathcal{M})_{X}
\qquad\text{resp.}\qquad
\mathrm{L}^{u}_{X/S}(\mathcal{M})\xrightarrow{\mathrm{L}(f)}
\mathrm{L}^{f\circ u}_{X'/S}(\mathcal{M})
\]
is a morphism of \(\OX\)-modules (\resp of \(\OS\)-modules).
\end{proposition}

This follows from Proposition~3.7 by functoriality in~\(\mathcal{M}\).

\begin{proposition}{3.8}\label{II.3.8}
Let \(X\) and \(Y\) be two functors over~\(S\). One has isomorphisms, functorial
in~\(\mathcal{M}\):
\begin{equation}
\mathrm{T}_{X/S}(\mathcal{M})\times_{S}\mathrm{T}_{Y/S}(\mathcal{M})
\isomto
\mathrm{T}_{(X\times_{S}Y)/S}(\mathcal{M}),
\tag{3.8.1}
\end{equation}
\begin{equation}
\mathrm{L}^{u}_{X/S}(\mathcal{M})\times_{S}\mathrm{L}^{v}_{Y/S}(\mathcal{M})
\isomto
\mathrm{L}^{(u,v)}_{(X\times_{S}Y)/S}(\mathcal{M}).
\tag{3.8.2}
\end{equation}
\end{proposition}

{\renewcommand{\thefootnote}{(28)}\nde{\normalfont What follows has been spelled out.}}
The first isomorphism follows from~1.2~(*), and the second is deduced from it
by the base change \((u,v)\colon S\to X\times_{S}Y\).

\begin{remark}{3.8.0}\label{II.3.8.0}
Let us observe that \textup{(3.8.1)} may also be interpreted as an isomorphism
of \(X\times_{S}Y\)-functors
\[
\bigl(\mathrm{T}_{X/S}(\mathcal{M})\times_{X}(X\times_{S}Y)\bigr)
\times_{X\times_{S}Y}
\bigl(\mathrm{T}_{Y/S}(\mathcal{M})\times_{Y}(X\times_{S}Y)\bigr)
\isomto
\mathrm{T}_{(X\times_{S}Y)/S}(\mathcal{M}).
\]
\end{remark}

\begin{corollary}{3.8.1}\label{II.3.8.1}
If \(X/S\) is equipped with an algebraic structure defined by finite cartesian
products, then \(\mathrm{T}_{X/S}(\mathcal{M})\) is equipped with a structure of
the same species, and the projection
\(\mathrm{T}_{X/S}(\mathcal{M})\to X\) is a morphism of this species of
structure.
\end{corollary}

\begin{proposition}{3.8. bis}\label{II.3.8.bis}
If \(X/S\) and \(Y/S\) satisfy \textup{(E)}, then \((X\times_{S}Y)/S\) satisfies
\textup{(E)}, and \textup{(3.8.1)} (resp.\ \textup{(3.8.2)}) is an isomorphism
of \(\mathcal{O}_{X\times_{S}Y}\)-modules (\resp of \(\OS\)-modules).
\end{proposition}

\begin{proof}%
{\renewcommand{\thefootnote}{(29)}\nde{\normalfont This proof has been added.}}
Suppose that \(X/S\) and \(Y/S\) satisfy \textup{(E)}. Then, by~3.5.1
and~\textup{(3.8.1)}, the same holds for \((X\times_{S}Y)/S\). Let us show
that \textup{(3.8.1)} is an isomorphism of
\(\mathcal{O}_{X\times_{S}Y}\)-modules.

Let \((x,y)\colon Z\to X\times_{S}Y\) be an \(S\)-morphism; taking~3.4.2 into
account, it suffices to see that the map
\[
\begin{aligned}
&\{\phi\in\Hom_{S}(\mathrm{I}_{Z}(\mathcal{M}),X)
\mid \phi\circ\varepsilon_{\mathcal{M}}=x\}
\times
\{\psi\in\Hom_{S}(\mathrm{I}_{Z}(\mathcal{M}),Y)
\mid \psi\circ\varepsilon_{\mathcal{M}}=y\}\\
&\qquad\longrightarrow
\{\theta\in\Hom_{S}(\mathrm{I}_{Z}(\mathcal{M}),X\times_{S}Y)
\mid \theta\circ\varepsilon_{\mathcal{M}}=(x,y)\}
\end{aligned}
\]
which to \((\phi,\psi)\) associates \(\phi\times\psi\) is a morphism of
\(\mathcal{O}(Z)\)-modules. But this is clear, for if
\(a\in\mathcal{O}(Z)\), then
\(a\cdot(\phi,\psi)=(\phi\circ a^{*},\psi\circ a^{*})\) is sent to
\[
(\phi\circ a^{*})\times(\psi\circ a^{*})
=(\phi\times\psi)\circ a^{*}
=a\cdot(\phi\times\psi).
\]
Likewise, using~3.2.1, one shows that \textup{(3.8.2)} is an isomorphism of
\(\OS\)-modules.
\end{proof}

\numpara{3.9.0}\label{II.3.9.0}%
{\renewcommand{\thefootnote}{(30)}\nde{\normalfont This paragraph has been added, in order to explain the introduction of the notion of \(S\)-H-functor.}}
If \(X\) is an \(S\)-group and if \(e\colon S\to X\) denotes its unit section,
one denotes:
\[
\Lie(X/S,\mathcal{M})=\mathrm{L}^{e}_{X/S}(\mathcal{M}),
\]
\ie \(\Lie(X/S,\mathcal{M})\) is defined by the cartesian square:
\[
\xymatrix{
\Lie(X/S,\mathcal{M})\ar[r]^{i}\ar[d]
& \mathrm{T}_{X/S}(\mathcal{M})\ar[d]^{p}\\
S\ar[r]^{e}
& X.
}
\]

By Corollary~3.8.1, the projection
\(p\colon\mathrm{T}_{X/S}(\mathcal{M})\to X\) is a morphism of \(S\)-groups, and
it follows that \(\Lie(X/S,\mathcal{M})\) is equipped with a structure of
\(S\)-group, and is isomorphic via~\(i\) to the kernel of~\(p\).

If, moreover, \(X/S\) satisfies condition~\textup{(E)}, one will see in
Proposition~3.9 that the structure of \(S\)-group on
\(\Lie(X/S,\mathcal{M})\), induced by that of~\(X\), \emph{coincides with the
structure of abelian group induced by the functoriality in~\(\mathcal{M}\)}
(\cf 3.5.1). In fact, this result is valid under the weaker hypothesis that
\(X\) be an \(S\)-functor in monoids or, more generally, an \(S\)-functor in
H-sets (\cf the definition below).

\begin{enoncerm}{Definitions}{3.9.0.1}\label{II.3.9.0.1}
a)~Let us introduce the following terminology:%
{\renewcommand{\thefootnote}{(31)}\nde{\normalfont This is inspired by the notion of H-space in topology.}}
an H-set is a set \(X\) equipped with a composition law with two-sided
unit, denoted \(e_{X}\) or simply~\(e\). If \(f\colon X\to Y\) is a morphism
of H-sets, its kernel \(\Ker f\) is \(f^{-1}(e_{Y})\); it is a
sub-H-set of~\(X\).

b)~An H-object in a category \(\mathcal{C}\) is defined in the usual way: it
is thus an object \(X\) of \(\mathcal{C}\), equipped with a morphism
\(X\times X\to X\) such that there exists a section of~\(X\) (over the final
object) possessing the properties of a two-sided unit. Every
\(\mathcal{C}\)-monoid, in particular every \(\mathcal{C}\)-group, is thus a
\(\mathcal{C}\)-H-object. In particular, an H-object of the category of
functors over the scheme~\(S\) will be called an \(S\)-H-functor.

c)~If \(X\) is an \(S\)-H-functor (for example, an \(S\)-group), and if
\(e\colon S\to X\) denotes the unit section of~\(X\), one denotes:
\[
\Lie(X/S,\mathcal{M})=\mathrm{L}^{e}_{X/S}(\mathcal{M})
\quad\text{and}\quad
\Lie(X/S)=\Lie(X/S,\OS).
\]
\end{enoncerm}

Let us then make explicit the following particular case of~3.8.1.%
{\renewcommand{\thefootnote}{(32)}\nde{\normalfont Corollary~3.9.0.2 has been added; it will be useful in~4.1.}}

\begin{corollary}{3.9.0.2}\label{II.3.9.0.2}
If \(X\) is an \(S\)-H-functor (\resp an \(S\)-group), then
\(\mathrm{T}_{X/S}(\mathcal{M})\) and \(\Lie(X/S,\mathcal{M})\) are also
\(S\)-H-functors (resp.\ \(S\)-groups), and one has morphisms of
\(S\)-H-functors (\resp of \(S\)-groups):
\[
\xymatrix{
\Lie(X/S,\mathcal{M})\ar[r]^{i}
& \mathrm{T}_{X/S}(\mathcal{M})\ar@<0.5ex>[r]^{p}
& X\ar@<0.5ex>[l]^{s}\,,
}
\]
where \(i\) is an isomorphism of \(\Lie(X/S)(\mathcal{M})\) onto \(\Ker p\) and
\(s\) is a section of~\(p\).
% typo?: printed Lie(X/S)(M), not Lie(X/S, M); kept as printed
\end{corollary}

\begin{proposition}{3.9}\label{II.3.9}
Let \(X\) be an \(S\)-H-functor satisfying \textup{(E)} with respect to~\(S\).
The structure of \(S\)-H-functor on \(\Lie(X/S,\mathcal{M})\) coming from that
of~\(X\) coincides with the structure of \(S\)-group defined in~3.5.1.
\end{proposition}

It follows from what has been said above that \(\Lie(X/S,\mathcal{M})\) is an
H-object in the category of \(\OS\)-modules. The proposition will then follow
from the following lemma:%
{\renewcommand{\thefootnote}{(33)}\nde{\normalfont Inspired by the standard proof that the fundamental group of an H-space is abelian.}}

\begin{lemma}{3.10}\label{II.3.10}
Let \(\mathcal{C}\) be a category. Let \(G\) be an H-object in the category of
\(\mathcal{C}\)-H-objects; \(G\) is thus a \(\mathcal{C}\)-H-object (whose
composition law we shall denote by \(f\colon G\times G\to G\)) equipped with
a morphism of \(\mathcal{C}\)-H-objects \(h\colon G\times G\to G\).%
{\renewcommand{\thefootnote}{(34)}\nde{\normalfont \(G\times G\) is equipped with the law \((G\times G)\times(G\times G)\to G\times G\), \(((x,z),(y,t))\mapsto(f(x,y),f(z,t))\).}}
Then \(f=h\) and \(f\) is commutative.
\end{lemma}

By taking the values of the functors on a variable argument, one reduces in
the usual way to checking the lemma when \(\mathcal{C}\) is the category of
sets. One thus has a set \(G\) and two maps \(f\), \(h\colon G\times G\to G\)
such that
\[
h\bigl(f(x,y),f(z,t)\bigr)=f\bigl(h(x,z),h(y,t)\bigr).
\]
On the other hand one has two elements of~\(G\), namely \(e\) and~\(u\), with
\(f(e,x)=f(x,e)=x\), \(h(u,x)=h(x,u)=x\). One sees first that
\[
h\bigl(f(u,y),f(x,u)\bigr)=f(x,y)=h\bigl(f(x,u),f(u,y)\bigr).
\]
In particular, for \(y=e\), resp.\ \(x=e\), one obtains, respectively,
\begin{align*}
x&=f(x,e)=h\bigl(f(u,e),f(x,u)\bigr)=h\bigl(u,f(x,u)\bigr)=f(x,u),\\
y&=f(e,y)=h\bigl(f(e,u),f(u,y)\bigr)=h\bigl(u,f(u,y)\bigr)=f(u,y),
\end{align*}
whence, substituting back into the original equality,
\[
h(y,x)=f(x,y)=h(x,y).
\]
This proves the lemma, and also Proposition~3.9. One then deduces from~3.9
the following corollaries.

\begin{corollary}{3.9.1}\label{II.3.9.1}
If \(X\) is an \(S\)-H-functor satisfying \textup{(E)} with respect to~\(S\),
every element of \(X(\mathrm{I}_{S}(\mathcal{M}))\) which projects onto the
unit element of \(X(S)\) is invertible.
\end{corollary}

\begin{corollary}{3.9.2}\label{II.3.9.2}
If \(X\) is an \(S\)-monoid satisfying \textup{(E)} with respect to~\(S\), an
element of \(X(\mathrm{I}_{S}(\mathcal{M}))\) is invertible if and only if its
image in \(X(S)\) is invertible.
\end{corollary}

\begin{corollary}{3.9.3}\label{II.3.9.3}
If \(X\) is an \(S\)-group satisfying \textup{(E)} with respect to~\(S\), the
two \(S\)-group laws on \(\Lie(X/S,\mathcal{M})\) coincide.
\end{corollary}

\begin{corollary}{3.9.4}\label{II.3.9.4}%
{\renewcommand{\thefootnote}{(35)}\nde{\normalfont This corollary has been added; it amplifies Remark~4.1.1.2 further on.}}
Let \(G\) be an \(S\)-group satisfying \textup{(E)} with respect to~\(S\). For
\(n\in\ZZ\), let \(n_{G}\colon G\to G\) be the morphism of \(S\)-functors
defined by \(g\mapsto g^{n}\). Then the derived morphism
\(\mathrm{L}(n_{G})\colon\Lie(G/S)\to\Lie(G/S)\) is the ``multiplication
by~\(n\)'', \ie the map which to every \(x\in\Lie(G/S)(S')\) associates
\(n\,x\).
\end{corollary}

Let us first observe that \(n_{G}\) is not in general a morphism of groups,
but that it preserves the unit section \(e\colon S\to G\), so the derived
morphism \(\mathrm{L}(n_{G})\) does indeed send
\(\Lie(G/S)=\mathrm{L}^{e}_{X/S}\) into itself.
% typo?: printed L^e_{X/S}, not L^e_{G/S}; kept as printed
If one denotes by~\(i\) the inclusion
\(\Lie(G/S)\hookrightarrow\mathrm{T}_{G/S}\), then \(\mathrm{L}(n_{G})\) is
defined by the equality \(i\bigl(\mathrm{L}(n_{G})(x)\bigr)=i(x)^{n}\), for
every \(S'\to S\) and every \(x\in\Lie(G/S)(S')\). Now, by~3.9, the two group
laws on \(\Lie(G/S)\) (the one coming from condition~\textup{(E)} and the one
coming from the law of~\(G\)) coincide, \ie one has
\(i(x)^{n}=i(n\,x)\), whence \(\mathrm{L}(n_{G})(x)=n\,x\).

Before drawing further consequences from Proposition~3.9, let us prove
another functoriality result:

\begin{proposition}{3.11}\label{II.3.11}
In the situation of Section~1, one has an isomorphism, functorial
in~\(\mathcal{M}\),
\[
\mathrm{T}_{\Hom_{Z/S}(X,Y)/S}(\mathcal{M})
\isomto
\Hom_{Z/S}\bigl(X,\mathrm{T}_{Y/Z}(\mathcal{M})\bigr).
\]
\end{proposition}

Indeed, one has by definition (\cf 3.1):
\[
\mathrm{T}_{\Hom_{Z/S}(X,Y)/S}(\mathcal{M})
=\Hom_{S}\bigl(\mathrm{I}_{S}(\mathcal{M}),\Hom_{Z/S}(X,Y)\bigr).
\]
By isomorphism~(4) of~1.1, applied to \(T=\mathrm{I}_{S}(\mathcal{M})\),
one has:
\[
\Hom_{S}\bigl(\mathrm{I}_{S}(\mathcal{M}),\Hom_{Z/S}(X,Y)\bigr)
\simeq
\Hom_{Z/S}\bigl(X,\Hom_{Z}(Z\times_{S}\mathrm{I}_{S}(\mathcal{M}),Y)\bigr).
\]
Taking account of the isomorphism
\(Z\times_{S}\mathrm{I}_{S}(\mathcal{M})\simeq\mathrm{I}_{Z}(\mathcal{M})\),
this gives
\begin{align*}
\mathrm{T}_{\Hom_{Z/S}(X,Y)/S}(\mathcal{M})
&\simeq
\Hom_{Z/S}\bigl(X,\Hom_{Z}(\mathrm{I}_{Z}(\mathcal{M}),Y)\bigr)\\
&=\Hom_{Z/S}\bigl(X,\mathrm{T}_{Y/Z}(\mathcal{M})\bigr).
\end{align*}

\begin{corollary}{3.11.1}\label{II.3.11.1}
If \(Y/Z\) satisfies \textup{(E)}, then \(\Hom_{Z/S}(X,Y)/S\) satisfies
\textup{(E)}, and the isomorphism of~3.11 respects the structures of
\(\mathcal{O}\)-modules over \(\Hom_{Z/S}(X,Y)\).%
{\renewcommand{\thefootnote}{(36)}\nde{\normalfont Setting \(H=\Hom_{Z/S}(X,Y)\), this implies, in particular, that \(\mathrm{T}_{H/S}(\mathcal{M})\) is equipped with a natural structure of module over \(\prod_{X\times_{S}H/H}\mathcal{O}_{H}\). This result seemed a little surprising to the editors; for this reason its proof has been spelled out.}}
\end{corollary}

\begin{proof}
Let \(\mathcal{M}\), \(\mathcal{N}\) be two free \(\OS\)-modules of finite
type. If \(Y/Z\) satisfies \textup{(E)}, then
\[
\mathrm{T}_{Y/Z}(\mathcal{M}\oplus\mathcal{N})
\simeq
\mathrm{T}_{Y/Z}(\mathcal{M})\times_{Y}\mathrm{T}_{Y/Z}(\mathcal{N}).
\]
The right-hand term is a subfunctor of
\(\mathrm{T}_{Y/Z}(\mathcal{M})\times_{S}\mathrm{T}_{Y/Z}(\mathcal{N})\), and
via the isomorphism of~1.2~(*), one obtains an isomorphism
\begin{align*}
&\Hom_{Z/S}\bigl(X,\mathrm{T}_{Y/Z}(\mathcal{M}\oplus\mathcal{N})\bigr)\\
&\qquad\simeq
\Hom_{Z/S}\bigl(X,\mathrm{T}_{Y/Z}(\mathcal{M})\bigr)
\times_{\Hom_{Z/S}(X,Y)}
\Hom_{Z/S}\bigl(X,\mathrm{T}_{Y/Z}(\mathcal{N})\bigr).
\end{align*}
Combined with~3.11, this yields:
\begin{align*}
&\mathrm{T}_{\Hom_{Z/S}(X,Y)/S}(\mathcal{M}\oplus\mathcal{N})\\
&\qquad\simeq
\mathrm{T}_{\Hom_{Z/S}(X,Y)/S}(\mathcal{M})
\times_{\Hom_{Z/S}(X,Y)}
\mathrm{T}_{\Hom_{Z/S}(X,Y)/S}(\mathcal{N}),
\end{align*}
so \(\Hom_{Z/S}(X,Y)\) satisfies \textup{(E)} with respect to~\(S\). This
proves the first assertion of the corollary.

Let us turn to the second. Write \(H=\Hom_{Z/S}(X,Y)\), and let us take an
\(S\)-morphism \(\Delta\colon H'\to\Hom_{Z/S}(X,Y)\), \ie a \(Z\)-morphism
\(\delta\colon H'\times_{S}X\to Y\), which thus makes \(H'\times_{S}X\) into
a \(Y\)-object.

On the one hand, one has a commutative diagram:
\[
\xymatrix@C=1.6em@R=2.2em{
\Hom_{H}\bigl(H',\Hom_{Z/S}(X,\mathrm{T}_{Y/Z}(\mathcal{M}))\bigr)
  \ar@{^{(}->}[r]\ar[d]
& \Hom_{S}\bigl(H',\Hom_{Z/S}(X,\mathrm{T}_{Y/Z}(\mathcal{M}))\bigr)\ar[d]\\
\Hom_{Y}\bigl(H'\times_{S}X,\mathrm{T}_{Y/Z}(\mathcal{M})\bigr)
  \ar@{^{(}->}[r]\ar[d]
& \Hom_{Z}\bigl(H'\times_{S}X,\mathrm{T}_{Y/Z}(\mathcal{M})\bigr)\ar[d]\\
\bigl\{\psi\in\Hom_{Z}(\mathrm{I}_{H'\times_{S}X}(\mathcal{M}),Y)
  \mid\psi\circ\varepsilon_{\mathcal{M}}=\delta\bigr\}
  \ar@{^{(}->}[r]
& \Hom_{Z}\bigl(\mathrm{I}_{H'\times_{S}X}(\mathcal{M}),Y\bigr).
}
\]
By~1.3, the action of \(\alpha\in\mathcal{O}(H'\times_{S}X)\) on
\(\Psi\in\Hom_{Y}(H'\times_{S}X,\mathrm{T}_{Y/Z}(\mathcal{M}))\) is given by:
for every \(U\to S\) and every
\((h,x)\in\Hom_{S}(U,H'\times_{S}X)\) (with \(U\) then lying over~\(Y\) via
\(\delta\circ(h,x)\)), one has:
\[
(\alpha\Psi)(h,x)=\alpha(h,x)\,\Psi(h,x),
\]
where \(\alpha(h,x)\in\mathcal{O}(U)\) acts on
\(\Psi(h,x)\in\mathrm{T}_{Y/Z}(\mathcal{M})(U)\) via the structure of
\(\mathcal{O}_{Y}\)-module on \(\mathrm{T}_{Y/Z}(\mathcal{M})\). By~3.4.2,
the latter is given, via the identification
\[
\begin{aligned}
&\Hom_{Y}\bigl(H'\times_{S}X,\mathrm{T}_{Y/Z}(\mathcal{M})\bigr)\\
&\quad=
\bigl\{\psi\in\Hom_{Z}(\mathrm{I}_{H'\times_{S}X}(\mathcal{M}),Y)
\mid\psi\circ\varepsilon_{\mathcal{M}}=\delta\bigr\},
\end{aligned}
\]
by: for every
\((m,h,x)\in\Hom_{S}(U,\mathrm{I}_{S}(\mathcal{M})\times_{S}H'\times_{S}X)\),
\begin{equation}
(\alpha\psi)(m,h,x)=\psi\bigl(m\cdot\alpha(h,x),h,x\bigr).
\tag{1}
\end{equation}

On the other hand, let us consider the tangent space
\(\mathrm{T}_{H/S}(\mathcal{M})=\Hom_{S}(\mathrm{I}_{S}(\mathcal{M}),H)\); one
has a commutative diagram:
\[
\xymatrix@C=1.6em@R=2.4em{
\Hom_{H}\bigl(H',\mathrm{T}_{H/S}(\mathcal{M})\bigr)
  \ar@{^{(}->}[r]\ar[d]
& \Hom_{S}\bigl(H',\mathrm{T}_{H/S}(\mathcal{M})\bigr)\ar[d]\\
\bigl\{\Phi\in\Hom_{S}(\mathrm{I}_{H'}(\mathcal{M}),H)
  \mid\Phi\circ\varepsilon_{\mathcal{M}}=\Delta\bigr\}
  \ar@{^{(}->}[r]\ar[d]_{(\ast)}
& \Hom_{S}\bigl(\mathrm{I}_{H'}(\mathcal{M}),H\bigr)\ar[d]\\
\bigl\{\phi\in\Hom_{Z}(\mathrm{I}_{H'\times_{S}X}(\mathcal{M}),Y)
  \mid\phi\circ\varepsilon_{\mathcal{M}}=\delta\bigr\}
  \ar@{^{(}->}[r]
& \Hom_{Z}\bigl(\mathrm{I}_{H'\times_{S}X}(\mathcal{M}),Y\bigr),
}
\]
where the bijection \((*)\) is given as follows (\cf 1.1~(2) and I~1.7.1):
for every \(U\to S\) and every
\((m,h,x)\in\Hom_{S}\bigl(U,\mathrm{I}_{S}(\mathcal{M})\times_{S}H'\times_{S}X\bigr)\)
(so that \(U\) lies over~\(Z\) via \(U\xrightarrow{x}X\to Z\)), one has
\(\Phi(m,h)\in\Hom_{Z}(X\times_{S}U,Y)\) and
\begin{equation}
\phi(m,h,x)=\Phi(m,h)\circ(x\times\mathrm{id}_{U})
\in\Hom_{Z}(U,Y).
\tag{\dag}
\end{equation}
By~3.4.2 (where one replaces \(X\) by \(\Hom_{Z/S}(X,Y)\) and \(X'\)
by~\(H'\)), the action of \(a\in\mathcal{O}(H')\) on
\(\Phi\in\Hom_{S}(\mathrm{I}_{H'}(\mathcal{M}),H)\) is given by: for every
\(U\to S\) and every
\((m,h)\in\Hom_{S}(U,\mathrm{I}_{S}(\mathcal{M})\times_{S}H')\),
\[
(a\Phi)(m,h)=\Phi\bigl(m\cdot a(h),h\bigr).
\]
Consequently, if \(\phi\) (resp.\ \(a\phi\)) is the element of
\(\Hom_{Z}(\mathrm{I}_{H'\times_{S}X}(\mathcal{M}),Y)\) associated with
\(\Phi\) (resp.\ \(a\Phi\)), one has, by~\((\dagger)\),
\begin{equation}
\begin{aligned}
(a\phi)(m,h,x)
&=\Phi\bigl(m\cdot a(h),h\bigr)\circ(x\times\mathrm{id}_{U})\\
&=\phi\bigl(m\cdot a(h),h,x\bigr).
\end{aligned}
\tag{2}
\end{equation}
Together with~(1), this shows that the isomorphism
\(\mathrm{T}_{H/S}(\mathcal{M})\isomto
\Hom_{Z/S}\bigl(X,\mathrm{T}_{Y/Z}(\mathcal{M})\bigr)\)
of~3.11.1 is an isomorphism of \(\mathcal{O}(H)\)-modules; moreover, for every
\(H'\to H\), the structure of \(\mathcal{O}(H')\)-module on
\(\Hom_{H}\bigl(H',\mathrm{T}_{H/S}(\mathcal{M})\bigr)\) extends, in a manner
functorial in~\(H'\), to a structure of
\(\mathcal{O}(H'\times_{S}X)\)-module.
\end{proof}

In particular, for \(Z=S\), one obtains the following corollary.

\begin{corollary}{3.11.2}\label{II.3.11.2}
One has an isomorphism, functorial in~\(\mathcal{M}\),
\[
\mathrm{T}_{\Hom_{S}(X,Y)/S}(\mathcal{M})
\isomto
\Hom_{S}\bigl(X,\mathrm{T}_{Y/S}(\mathcal{M})\bigr).
\]
Moreover, if \(Y/S\) satisfies \textup{(E)}, then \(\Hom_{S}(X,Y)/S\)
satisfies \textup{(E)}, and the preceding isomorphism respects the structures
of \(\mathcal{O}\)-modules over \(\Hom_{S}(X,Y)\).
\end{corollary}

{\renewcommand{\thefootnote}{(37)}\nde{\normalfont The following sentences have been added.}}
Let \(u\colon X\to Y\) be an \(S\)-morphism; one identifies it with the
constant morphism \(u\colon S\to\Hom_{S}(X,Y)\) such that \(u(f)=u\) for
every \(f\colon S'\to S\). One then sees at once that the fibered product of
\(u\) and of
\(\Hom_{S}\bigl(X,\mathrm{T}_{Y/S}(\mathcal{M})\bigr)\to\Hom_{S}(X,Y)\)
is identified with
\(\Hom_{Y/S}\bigl(X,\mathrm{T}_{Y/S}(\mathcal{M})\bigr)\), where \(X\) lies
over~\(Y\) via~\(u\). Consequently, from the definition of
\(\mathrm{L}^{u}_{\Hom_{S}(X,Y)/S}(\mathcal{M})\) and from the preceding
corollary, one deduces the following:
